\documentclass[14pt, a4paper]{article}
\usepackage[fontsize=14pt]{fontsize}

\setcounter{page}{1}

\usepackage{graphicx} % Required for inserting images
\graphicspath{{../../results/}}

% Для кириллицы и русского языка
\usepackage[utf8]{inputenc}
\usepackage[english,russian]{babel}

% Заголовок содержания по ГОСТ: заглавными (babel может переопределять captions)
\addto\captionsrussian{\renewcommand{\contentsname}{СОДЕРЖАНИЕ}}

% Полуторный интервал
\usepackage{setspace}
\onehalfspacing

% Поля
\usepackage[left=30mm, right=15mm, top=20mm, bottom=20mm]{geometry}

% Отступы в 1,25см
\setlength{\parindent}{1.25cm}

% Применять отступы на первом абзаце в разделах и подразделах
\usepackage{indentfirst}

% Нумерация в нижней части включена по умолчанию

% Форматирование заголовков
\usepackage{titlesec}

% Разделы
\titleformat{\section}
{}
{\hspace{1.25cm}\arabic{section}}
{1ex}{}

\titleformat{\subsection}
{}
{\hspace{1.25cm}\arabic{section}.\arabic{subsection}}
{1ex}{}

\titleformat{\subsubsection}
{}
{\hspace{1.25cm}\arabic{section}.\arabic{subsection}.\arabic{subsubsection}}
{1ex}{}

% Структурные элементы
\titleformat{\paragraph}
{\centering}
{\hspace{1.25cm}\arabic{section}}
{1ex}{}

\newcommand{\structheader}[1]{%
	\phantomsection%
	\clearpage
	\paragraph*{{#1}}%
	\addcontentsline{toc}{section}{#1}%
}

% Заголовок, который должен попасть в содержание отчёта
\newcommand{\reportheader}[1]{%
	\phantomsection%
	\clearpage
	\paragraph*{{#1}}%
	\addcontentsline{toc}{section}{#1}%
}

\usepackage{fontspec}
\setmainfont{Times New Roman}


% Возможность вставки изображений и таблиц в нужное место
\usepackage{float}

% Для вставки PDF страниц
\usepackage{pdfpages}

% Таблицы с автоматическим выравниванием по ширине
\usepackage{array}
\usepackage{tabularx}

% Тире вместо двоеточия для подписей иллюстраций
\usepackage[figurename=Рисунок]{caption}
\DeclareCaptionLabelSeparator{tirer}{ -- }
\captionsetup{labelsep=tirer}

% Для форматирования списков (например, условных обозначений)
\usepackage{enumitem}

% ГОСТ-подобные абзацные отступы в списках:
% сам список начинается с абзацного отступа (\parindent).
% Для itemize используем маркер в виде «палки» (тире), а не точки.
\setlist[itemize]{leftmargin=\parindent, labelindent=0pt, listparindent=\parindent, label=--}
\setlist[enumerate]{leftmargin=\parindent, labelindent=0pt, listparindent=\parindent, label=\arabic*)}

% Для форматирования содержания
\usepackage{tocloft}

\setlength{\cftbeforesecskip}{0pt}
\renewcommand{\cfttoctitlefont}{\normalfont\fontsize{14pt}{14pt}\selectfont}
\renewcommand{\cftaftertoctitle}{}
\renewcommand{\cftbeforetoctitleskip}{-1em}
\renewcommand{\cftsecpresnum}{}
\newlength\mylength
\settowidth\mylength{\cftsecpresnum}
\addtolength\cftsecnumwidth{\mylength}

\renewcommand{\cftbeforetoctitleskip}{-1em}
\renewcommand{\cftsecfont}{\normalfont}
\renewcommand{\cftsubsecfont}{\normalfont}
\renewcommand{\cftsubsubsecfont}{\normalfont}
\renewcommand{\cftsecleader}{\cftdotfill{\cftdotsep}}
\renewcommand{\cftsecpagefont}{\normalfont}

\usepackage{listings}
\usepackage{xcolor}
\usepackage{ragged2e}

\lstdefinestyle{mystyle}{
	basicstyle=\ttfamily\small,
	breaklines=true,                 % Enable line wrapping
	frame=single,                     % Add a single frame around the code
	framesep=5pt,                     % Padding between frame and code
	rulecolor=\color{black},           % Frame color
	backgroundcolor=\color{white},     % Background color
	postbreak=\mbox{\textcolor{red}{$\hookrightarrow$}\space}, % Mark line breaks
	showstringspaces=false,           % Don't show spaces in strings
	tabsize=2,                        % Set tab size
	captionpos=b,                      % Caption position (bottom)
	columns=flexible
}

\lstset{style=mystyle}

% Настройки переносов: разрешаем перенос по слогам (в т.ч. русских слов)
% и при этом избегаем грубого \sloppy.
\lefthyphenmin=2
\righthyphenmin=2
\hyphenpenalty=300
\exhyphenpenalty=50
\tolerance=2000
\emergencystretch=1.5em

\usepackage[hidelinks]{hyperref}

\newcommand{\srcref}[1]{\hyperlink{src:i#1}{[#1]}}
\newcommand{\srcrange}[2]{\srcref{#1}--\srcref{#2}}

% Ссылки на список источников ОСНОВНОГО ОТЧЁТА (перед Приложением 1)
\newcommand{\repsrcref}[1]{\hyperlink{rep:i#1}{[#1]}}
\newcommand{\repsrcrange}[2]{\repsrcref{#1}--\repsrcref{#2}}

% Многоэлементные ссылочные конструкции:~srcmulti{8,11,10} -> [1; 2; 3]
\usepackage{xparse}
\ExplSyntaxOn
\seq_new:N \l__src_items_seq
\NewDocumentCommand{\srcmulti}{m}
{
	\seq_set_split:Nnn \l_tmpa_seq { , } { #1 }
	\seq_clear:N \l__src_items_seq
	\seq_map_inline:Nn \l_tmpa_seq { \seq_put_right:Nn \l__src_items_seq { \hyperlink{src:i##1}{##1} } }
	[\seq_use:Nn \l__src_items_seq {;~}]
}

\seq_new:N \l__repsrc_items_seq
\NewDocumentCommand{\repsrcmulti}{m}
{
	\seq_set_split:Nnn \l_tmpa_seq { , } { #1 }
	\seq_clear:N \l__repsrc_items_seq
	\seq_map_inline:Nn \l_tmpa_seq { \seq_put_right:Nn \l__repsrc_items_seq { \hyperlink{rep:i##1}{##1} } }
	[\seq_use:Nn \l__repsrc_items_seq {;~}]
}
\ExplSyntaxOff


\begin{document}

	% 1) Титульный лист
	
\includepdf[pages=1]{title.pdf}

	\tableofcontents
	\clearpage
	\reportheader{ВВЕДЕНИЕ}

	Рост объёмов данных в современных поисковых системах приводит к необходимости разбиения поискового индекса на части (шарды), размещаемые на разных узлах кластера.
	При обработке запроса необходимо определить, какие шарды могут содержать релевантные документы, и направить запрос только в них.
	Число опрашиваемых шардов (fanout) напрямую влияет на время отклика и нагрузку на систему.

	Актуальность исследования обусловлена тем, что классический подход к шардированию~--- хеш-разбиение терминов~--- распределяет термины случайно, не учитывая их семантические связи.
	Это приводит к высокому fanout, поскольку запрос, содержащий несколько связанных терминов, направляется в несколько разных шардов.
	Семантический шардинг, при котором совместно встречающиеся термины группируются в одном шарде, позволяет снизить fanout и, следовательно, уменьшить время обработки запроса.

	Цель работы~--- исследование производительности поисковых индексов, построенных на основе кластеризации поисковых запросов.

	Для достижения поставленной цели необходимо выполнить следующие задачи.
	\begin{enumerate}
		\item осуществить поиск существующих решений и подходов к шардированию в задаче построения поискового индекса;
		\item выполнить поиск набора данных поисковых запросов для кластеризации с целью шардирования индекса на основе полученных кластеров;
		\item реализовать прототип поискового индекса на основе полученных кластеров для найденного набора данных;
		\item исследовать производительность реализованного решения в сравнении с классическими способами построения индекса.
	\end{enumerate}

	\begin{table}[H]
	\centering
	\caption{Индивидуальное задание на практику}
	\label{tab:individual_task}
	\begin{tabularx}{\textwidth}{|>{\raggedright\arraybackslash}p{0.9cm}|>{\raggedright\arraybackslash}p{3cm}|>{\raggedright\arraybackslash}X|}
	\hline
	\textbf{№ этапа} & \textbf{Название этапа} & \textbf{Задание} \\
	\hline
	1 & Инструктаж обучающегося & Инструктаж обучающегося по ознакомлению с требованиями охраны труда, техники безопасности, пожарной безопасности, а также правилами внутреннего трудового распорядка \\
	\hline
	2 & Анализ существующих решений & Осуществить поиск существующих решений и подходов к шардированию в задаче построения поискового индекса \\
	\hline
	3 & Поиск датасета поисковых запросов для кластеризации & Выполнить поиск датасета поисковых запросов для кластеризации с целью шардирования индекса на основе полученных кластеров \\
	\hline
	4 & Реализация прототипа поискового индекса на основе кластеров & Реализовать прототип поискового индекса на основе полученных кластеров для найденного датасета \\
	\hline
	5 & Исследовать производительность реализованного решения & Исследовать производительность реализованного решения в сравнении с классическими способами построения индекса \\
	\hline
	\end{tabularx}
	\end{table}

	\begin{table}[H]
	\centering
	\caption*{Продолжение таблицы~\ref{tab:individual_task}}
	\begin{tabularx}{\textwidth}{|>{\raggedright\arraybackslash}p{0.9cm}|>{\raggedright\arraybackslash}p{3cm}|>{\raggedright\arraybackslash}X|}
	\hline
	\textbf{№ этапа} & \textbf{Название этапа} & \textbf{Задание} \\
	\hline	
	6 & Оформление отчёта, получение отзыва с оценкой от руководителя & Написание отчёта, оформление в соответствии с требованиями, получение отзыва от руководителей. Оформление отчётных документов в соответствии с требованиями: 1. Оформление отчёта должно быть выполнено в соответствии с методическим пособием (https://books.ifmo.ru/file/pdf/2622.pdf, разделы "Набор текста", "Нумерация страниц и разделов", "Заголовки разделов", "Оформление рисунков", "Оформление таблиц", "Оформление формул", "Оформление ссылок на источники", "Оформление списка использованных источников", "Оформление приложений"). 2. Структура документа: титульный лист, введение, основная часть, заключение, приложения. 3. В основной части подробно описывается выполнение задач всех этапов. Получить отзыв с оценкой у руководителя практики в модуле Практика. \\
	\hline
	\end{tabularx}
	\end{table}

	\section{Анализ существующих решений}

	\subsection{Постановка задачи}

	Инвертированный индекс~--- основная структура данных в информационном поиске, отображающая термин в список документов, содержащих этот термин.
	При поиске по запросу система извлекает списки документов для всех терминов запроса, пересекает или объединяет их и ранжирует результаты (например, по BM25).

	При росте объёмов данных инвертированный индекс перестаёт помещаться на одном сервере и должен быть разбит на части~--- шарды.
	Возникает задача \textit{маршрутизации запроса}: определить, какие шарды могут содержать релевантные документы, и направить запрос только в них, минимизируя число опрашиваемых шардов при сохранении полноты поиска (recall).

	\subsection{Подходы к шардингу поискового индекса}

	Существует два основных подхода к шардированию инвертированного индекса~\srcmulti{1,2,3}.

	1. \textit{Шардирование по документам}.
	Каждый документ назначается в шард по $\mathrm{hash}(\mathrm{doc\_id}) \bmod N$.
	Это стандартный подход в Elasticsearch и большинстве распределённых поисковых систем.
	Документы распределены равномерно, однако каждый шард содержит все термины, и для обработки запроса необходимо опрашивать \textit{все} шарды (fanout~$= N$).

	2. \textit{Шардирование по терминам}.
	Каждый термин назначается в определённый шард, и шард содержит списки документов только для своих терминов.
	Документ назначается во все шарды, где есть его термины.
	Запрос маршрутизируется только в шарды, содержащие термины запроса.
	Fanout зависит от числа уникальных шардов терминов запроса.

	Однако оба подхода имеют свои недостатки, для решения которых предлагается такой подход, как семантический шардинг (по кластерам связанных терминов).
	Он отличается от шардирования по терминам тем, что назначение терминов в шарды производится не случайно, а на основе кластеризации, то есть совместно встречающиеся термины группируются в один шард.
	Это позволяет уменьшить fanout, поскольку связанные термины запроса с высокой вероятностью оказываются в одном шарде (особенно эффективно для часто встречающихся вместе терминов).

	\subsection{Обзор релевантных работ}

	Задача разбиения поискового индекса и маршрутизации запросов рассматривается в фундаментальной работе Маннинга и др., где описаны стратегии разбиения по документам и по терминам, а также методы оценки полноты маршрутизации~\srcref{4}.
	Ранний подход к определению релевантных коллекций для поиска представлен в системе GIOSS, где для каждого источника строилась статистическая модель на основе частотностей терминов, позволявшая оценить вероятность наличия релевантных документов без прямого обращения к коллекции~\srcref{5}.

	Метод Taily предлагает подход к выбору шардов на основе распределения оценок терминов в каждом шарде~\srcref{6}.
	Для каждого шарда строится приближение распределения оценок (tail распределение), и шарды ранжируются по вероятности содержать высокооценённые документы.
	В отличие от GIOSS, Taily не требует хранения глобальной статистики по всем терминам и лучше работает при большом числе шардов.

	Система BitFunnel реализует шардирование по терминам в промышленном поисковике Bing~\srcref{7}.
	Вместо традиционного инвертированного индекса BitFunnel использует битовую матрицу <<термин--шард>>, где каждый бит указывает наличие термина в шарде, что обеспечивает компактное представление и быструю маршрутизацию.
	Документы назначаются во все шарды, содержащие их термины, а запрос направляется только в шарды с совпадающими терминами.
	При масштабировании до миллиардов документов основной проблемой становится рост числа шардов и увеличение fanout.
	Система Maguro также реализует шардирование по терминам~\srcref{8}.

	В отличие от перечисленных подходов работы по эмбеддинговому поиску предлагают использовать векторные представления запросов и документов для маршрутизации~\srcref{9}.
	Однако такие подходы требуют значительных вычислительных ресурсов на обучение моделей и построение эмбеддингов.

	В текущей работе рассматривается промежуточный подход, который находится между статистическим и нейросетевым.
	Вместо обучения моделей для векторных представлений используется кластеризация на графе совстречаемости терминов, что позволяет группировать семантически связанные термины без обучения нейронной сети~\srcref{10}.

	\clearpage
	\section{Выбор набора данных}

	\subsection{Требования к данным}

	Для проведения исследования необходим набор данных, удовлетворяющий следующим требованиям:
	\begin{itemize}
		\item наличие поисковых запросов и документов для них;
		\item достаточный объём для оценки шардинга (более 100\,000 документов);
		\item публичная доступность для воспроизводимости исследования.
	\end{itemize}

	\subsection{Обзор кандидатов}

	Рассмотрены следующие наборы данных:
	\begin{enumerate}
		\item \textit{MS~MARCO v1}~--- коллекция из 8.8 млн документов и 1 млн запросов, один из стандартных тестовых наборов для задач информационного поиска;
		\item \textit{MS~MARCO v2.1}~--- обновлённая версия с улучшенной разметкой, содержит 808\,731 запросов в обучающем разделе, каждому из которых соответствует 10 документов~\srcref{13};
		\item \textit{TREC CAR}~--- коллекция для ответа на вопросы по статьям Википедии;
		\item \textit{Natural Questions}~--- набор данных от Google с вопросами и ответами из Википедии, формат вопрос--ответ.
	\end{enumerate}

	\subsection{MS~MARCO v2.1}

	Для исследования выбран набор данных MS~MARCO v2.1 по следующим причинам:
	\begin{enumerate}
		\item большой объём (808\,731 запросов с 10 документами каждый), обеспечивающий достаточную статистику для оценки шардинга;
		\item каждый пример содержит текст запроса, текст документа и признак \texttt{is\_selected}, что позволяет построить эталонные данные для оценки полноты маршрутизации;
		\item публичная доступность через библиотеку HuggingFace \texttt{datasets};
		\item широкое использование в исследованиях по информационному поиску, обеспечивающее сравнимость результатов.
	\end{enumerate}

	\section{Реализация прототипа}

	\subsection{Построение индекса}

	Для построения индекса необходимо выполнить три основных этапа:
	\begin{enumerate}
		\item построить граф совстречаемости терминов с весами NPMI;
		\item кластеризовать граф алгоритмом Leiden с получением отображения \textit{термин~$\to$~шард};
		\item назначить каждый документ во все шарды, содержащие его термины.
	\end{enumerate}

	\subsection{Построение NPMI-графа совстречаемости}

	Для построения графа совстречаемости каждый документ представляется бинарным вектором терминов (факт встречаемости без учёта частоты) с фильтрацией стоп-слов и ограничениями на частотность.
	По матрице встречаемости вычисляется матрица совстречаемости (термины считаются совстретившимися, если они входят в один и тот же запрос), после чего для каждой пары терминов с достаточным числом совместных упоминаний рассчитывается нормированная взаимная информация (NPMI):
	\begin{equation}
		\mathrm{NPMI}(x,y) = \frac{\log\frac{P(x,y)}{P(x)P(y)}}{-\log P(x,y)}
	\end{equation}
	NPMI нормирует PMI к отрезку $[-1, 1]$, что удобно для сравнения ассоциаций разной частотности.
	Положительное значение NPMI указывает на то, что термины встречаются совместно чаще, чем при случайном распределении~\srcref{10}.
	Пары с отрицательным NPMI отбрасываются.
	В результате формируется взвешенный граф, где вершины~--- термины, рёбра~--- значимые совместные встречаемости, веса~--- значения NPMI.
	На наборе данных MS~MARCO v2.1 (train) граф содержит 15\,000 терминов (минимальное количество совстречаемости двух терминов для принятия~--- 2, минимальная частота появления в документах~--- 2, максимальная доля документов, содержащих термин~--- 0.5).

	\subsection{Кластеризация алгоритмом Leiden}

	Разбиение терминов на кластеры выполняется с помощью алгоритма Leiden на построенном NPMI-графе~\srcmulti{11,12}.
	На наборе данных MS~MARCO v2.1 (train) алгоритм выделил 253 кластера (для разрешения $\gamma = 1.0$).

	Термины, не вошедшие в граф совстречаемости (низкочастотные или слишком частотные), назначаются в шарды по хеш-функции.
	В текстах документов содержится 139\,830 уникальных терминов, из них 124\,830 не попали в граф (построенный по запросам) и были назначены по хешу в дополнительные шарды.

	На рисунках~\ref{fig:npmi-clusters}--\ref{fig:tsne} представлены граф, распределение размеров, облака слов и t-SNE визуализация полученных кластеров соответственно.
	\begin{figure}[H]
		\centering
		\includegraphics[width=\textwidth]{npmi_graph_clusters.pdf}
		\caption{Визуализация кластеров NPMI-графа (топ-5 сообществ)}
		\label{fig:npmi-clusters}
	\end{figure}

	\begin{figure}[H]
		\centering
		\includegraphics[width=\textwidth]{cluster_size_distribution.pdf}
		\caption{Распределение размеров кластеров}
		\label{fig:cluster-sizes}
	\end{figure}

	\begin{figure}[H]
		\centering
		\includegraphics[width=\textwidth]{word_clouds_top25.pdf}
		\caption{Облака слов для топ-25 кластеров}
		\label{fig:wordclouds}
	\end{figure}

	\begin{figure}[H]
		\centering
		\includegraphics[width=\textwidth]{tsne_2d_projection.pdf}
		\caption{Проекция t-SNE терминов}
		\label{fig:tsne}
	\end{figure}

	Облака слов (рисунок~\ref{fig:wordclouds}) показывают, что каждый кластер образует тематически связную группу терминов.
	Проекция t-SNE (рисунок~\ref{fig:tsne}) демонстрирует чёткую локализацию кластеров в двумерном пространстве.

	\subsection{Назначение документов в шарды}

	Каждый документ назначается во все шарды, содержащие хотя бы один его термин.
	Это гарантирует полноту маршрутизации при поиске, но приводит к дублированию документов между шардами.
	Степень дублирования зависит от стратегии разбиения.
	При семантической кластеризации связанные термины оказываются в одном шарде, что приводит к меньшему дублированию по сравнению с хеш-разбиением.

	\section{Исследование производительности}

	\subsection{Методика исследования}

	Сравниваются две стратегии разбиения терминов по шардам:
	\begin{enumerate}
		\item \textit{Hash}~--- хеш-разбиение (по терминам);
		\item \textit{Leiden}~--- семантическая кластеризация.
	\end{enumerate}
	Обе стратегии назначают каждый документ во все шарды, содержащие его термины, и маршрутизируют запрос во все шарды, содержащие термины запроса.
	Набор данных MS~MARCO v2.1, раздел train (извлечено 120\,624 уникальных запроса и 129\,020 документов);
	В качестве метрик используются:
	\begin{enumerate}
		\item \textit{fanout}~--- число шардов, в которые направляется запрос;
		\item \textit{дублирование}~--- среднее число шардов, в которые попадает один документ;
		\item \textit{избыточность хранения}~--- отношение суммарного числа документов во всех шардах к числу уникальных документов.
	\end{enumerate}

	\subsection{Результаты исследования}

	В таблице~\ref{tab:partition-stats} приведены основные характеристики построенных разбиений.
	\begin{table}[H]
	\centering
	\caption{Статистика разбиения по шардам}
	\label{tab:partition-stats}
	\begin{tabularx}{\textwidth}{|>{\raggedright\arraybackslash}X|>{\raggedright\arraybackslash}X|>{\raggedright\arraybackslash}X|}
	\hline
	\textbf{Показатель} & \textbf{Hash} & \textbf{Leiden} \\
	\hline
	Дублирование & 24.4× & 15.1× \\
	\hline
	Избыточность хранения & 36.1× & 23.7× \\
	\hline
	\end{tabularx}
	\end{table}
	Дублирование документов для Leiden индекса на 38\% меньше, чем при хеш-шардировании (15.1× против 24.4×), избыточность хранения~--- на 34\% меньше (23.7× против 36.1×).

	На рисунке~\ref{fig:fanout} показано распределение fanout для 100\,000 запросов.
	\begin{figure}[H]
		\centering
		\includegraphics[width=\textwidth]{fanout_distribution.pdf}
		\caption{Распределение fanout}
		\label{fig:fanout}
	\end{figure}

	В таблице~\ref{tab:fanout} приведены метрики для fanout.
	\begin{table}[H]
	\centering
	\caption{Метрики fanout (для 100\,000 запросов)}
	\label{tab:fanout}
	\begin{tabularx}{0.7\textwidth}{|>{\raggedright\arraybackslash}X|>{\raggedright\arraybackslash}X|>{\raggedright\arraybackslash}X|}
	\hline
	\textbf{Метрика} & \textbf{Hash} & \textbf{Leiden} \\
	\hline
	Среднее & 3.21 & 2.29 \\
	\hline
	Медиана & 3 & 2 \\
	\hline
	P95 & 5 & 4 \\
	\hline
	\end{tabularx}
	\end{table}
	Средний fanout для Leiden индекса на 29\% ниже, чем при хеш-шардировании (2.29 против 3.21).

	Leiden достигает той же полноты маршрутизации при меньшем fanout, что следует из рисунка~\ref{fig:recall-fanout}, на котором представлена зависимость routing recall от количества опрашиваемых шардов.
	Такой выигрыш наблюдается из-за того, что предлагаемый подход учитывает семантическую связность кластеров.
	\begin{figure}[H]
		\centering
		\includegraphics[width=\textwidth]{recall_vs_fanout.pdf}
		\caption{Зависимость routing recall от числа опрашиваемых шардов}
		\label{fig:recall-fanout}
	\end{figure}

	\subsection{Ограничения}

	Наряду с преимуществами выявлены следующие ограничения:
	\begin{enumerate}
		\item дублирование остается достаточно большим (15.1×). Дальнейшее снижение дублирования приведет к уменьшению полноты маршрутизации;
		\item результат кластеризации существенно зависит от токенизации и фильтрации стоп-слов;
		\item дрифт данных требует полного перестроения всех кластеров, а значит и индексов;
		\item параметры построения графа (пороги частотности) и кластеризации (разрешение) влияют на число и размер кластеров;
		\item построение индексов для большого набора данных (больше миллиона документов) может занять достаточно большое количество времени.
	\end{enumerate}

	\structheader{ЗАКЛЮЧЕНИЕ}

	По результатам проведённого исследования можно сделать следующие выводы:
	\begin{enumerate}
		\item выполнен анализ существующих подходов к шардированию поискового индекса;
		\item выбран набор данных MS~MARCO v2.1 (извлечена подвыборка из 120\,624 запросов и 129\,020 документов), обеспечивающий достаточный объём данных и публичную доступность для воспроизводимости исследования;
		\item реализован прототип семантического шардированного индекса на основе кластеризации запросов;
		\item проведены сравнительные исследования с хеш-шардированием (по терминам): средний fanout при Leiden-шардировании на 29\% ниже (2.3 против 3.2), дублирование документов на 38\% меньше (15.1× против 24.4×), routing recall при полном fanout составляет 100\% для обеих стратегий, при ограниченном fanout Leiden обеспечивает более высокий recall, чем Hash;
		\item выявлены ограничения: высокое дублирование (15.1×), увеличение числа шардов из-за хеш-назначения терминов вне графа, более высокая латентность шардированного поиска по сравнению с полным индексом при отсутствии распределения.
	\end{enumerate}

	Перспективными направлениями дальнейших исследований являются:
	\begin{enumerate}
		\item кластеризация с учетом содержимого документов;
		\item инкрементальное обновление кластеров при потоковом поступлении данных;
		\item селективное назначение документов в шарды с ограничением дублирования;
		\item оценка влияния кластеризации на итоговые метрики качества поиска (NDCG, MRR);
		\item оптимизация параметров (пороги частотности, разрешение Leiden) для различных сценариев использования.
	\end{enumerate}

	\structheader{СПИСОК ИСПОЛЬЗОВАННЫХ ИСТОЧНИКОВ}

	\begin{enumerate}[wide=0pt, leftmargin=0pt, labelindent=\parindent, listparindent=\parindent, nosep, label=\arabic*]
		\item \hypertarget{src:i1}{} Fahad A., Alshatri N., Tari Z., Alamri A., Khalil I., Zomaya A. Y., Foufou S., Bouras A. A survey of clustering algorithms for big data: Taxonomy and empirical analysis // IEEE Transactions on Emerging Topics in Computing. -- 2014. -- Т. 2, № 3. -- С. 267--279.
		\item \hypertarget{src:i2}{} Wani A. A. Comprehensive analysis of clustering algorithms: exploring limitations and innovative solutions // PeerJ Computer Science. -- 2024. -- Т. 10. -- e2286.
		\item \hypertarget{src:i3}{} Wegmann M., Zipperling D., Hillenbrand J., Fleischer J. A review of systematic selection of clustering algorithms and their evaluation // arXiv preprint arXiv:2106.12792. -- 2021.
		\item \hypertarget{src:i4}{} Manning C. D., Raghavan P., Sch\"utze H. Introduction to Information Retrieval. -- Cambridge University Press, 2008. -- DOI: 10.1017/cbo9780511809071.
		\item \hypertarget{src:i5}{} Gravano L., Garc\'ia-Molina H., Tomasic A. GIOSS: Text-source discovery over the Internet // ACM Transactions on Database Systems. -- 1999. -- Т. 24, № 2. -- С. 229--264. DOI: 10.1145/320248.320252.
		\item \hypertarget{src:i6}{} Aly R., Hiemstra D., Demeester T. Taily: shard selection using the tail of score distributions // Proceedings of the 36th International ACM SIGIR Conference on Research and Development in Information Retrieval. -- 2013. -- С. 673--682. DOI: 10.1145/2484028.2484033.
		\item \hypertarget{src:i7}{} Goodwin M., Hopcroft M., Luu D., Clemmer A. V., Chieu H., Elnikety S. BitFunnel: a space-optimal signature index for search // Proceedings of the 40th International ACM SIGIR Conference on Research and Development in Information Retrieval. -- 2017. -- С. 593--602. DOI: 10.1145/3077136.3080789.
		\item \hypertarget{src:i8}{} Yamamoto Y., Shirahama K., Uehara K. Maguro: A Distributed Full-Text Search Engine with Term Partitioning // Proceedings of the 22nd International Conference on Extending Database Technology (EDBT). -- 2019. -- С. 546--549.
		\item \hypertarget{src:i9}{} Huang J.-T., Sharma A., Sun S., Xia L., Zhang D., Pronin P., Padmanabhan S., Ottaviano G., Yang L. Embedding-based Retrieval in FB Search // Proceedings of the 26th ACM SIGKDD International Conference on Knowledge Discovery and Data Mining. -- 2020. -- С. 2553--2561. DOI: 10.1145/3394486.3403305.
		\item \hypertarget{src:i10}{} Williams C. K. I. On Suspicious Coincidences and Pointwise Mutual Information // Neural Computation. -- 2022. -- Т. 34, № 10. -- С. 2037--2046. DOI: 10.1162/neco\_a\_01533.
		\item \hypertarget{src:i11}{} Traag V. A., Waltman L., van Eck N. J. From Louvain to Leiden: guaranteeing well-connected communities // Scientific Reports. -- 2019. -- Т. 9, № 1. -- С. 5233. DOI: 10.1038/s41598-019-41695-z.
		\item \hypertarget{src:i12}{} Chunaev P. Community detection in node-attributed social networks: A survey // Computer Science Review. -- 2020. -- Т. 37. -- Article 100286. DOI: 10.1016/j.cosrev.2020.100286.
		\item \hypertarget{src:i13}{} Craswell N., Mitra B., Yilmaz E., Campos D., Voorhees E. M. Overview of the TREC 2021 Deep Learning Track // Text REtrieval Conference (TREC). -- 2021.
	\end{enumerate}

	\end{document}
